\section{Introduction}
\label{sec:introduction}

Surface electromyography (sEMG) pattern recognition is the dominant approach to multifunction myoelectric control~\cite{hudgins1993new,englehart2003robust,scheme2011electromyogram,farina2014extraction}. Recognition generally improves with more recording sites, but every electrode adds cost, wiring, contact area and a point of failure in a wearable device~\cite{vujaklija2016new}. This makes it attractive to identify a reduced set of electrodes that preserves recognition performance~\cite{phinyomark2012feature,liu2014emg,altimemy2013classification,stango2015spatial,qu2021reduce}.

Three methodological problems limit what published channel-reduction results can tell a device designer. First, most reductions are selected and evaluated within the same users. \rev{User-dependent evaluation is appropriate for personalised systems, but it answers a different question from the one a device intended for many users must answer: a layout chosen for one person need not transfer to the next, and subject-independent recognition is generally harder than user-dependent recognition~\cite{phinyomark2018emg,cote2019deep}. Separately, if data from a test user influences scaling, ranking or tuning, the reported subset is optimistically biased.} Second, feature-attribution methods such as Shapley additive explanations~\cite{lundberg2017unified,lundberg2020local} are increasingly used to rank channels~\cite{salmi2025shapaad}, but attributions can be unstable across data splits~\cite{visani2022statistical,kelodjou2023shaping,xiang2026roshap}, and a ranking computed on the evaluation data is itself a form of leakage. Third, the chosen number of channels is usually read off an accuracy curve by eye, with no pre-stated tolerance and no statement of what ``no meaningful loss'' means~\cite{schuirmann1987comparison,walker2011understanding}. A margin chosen after the curve is seen cannot support a non-inferiority claim~\cite{campbell2018noninferiority}.

This study tests explainability-guided channel reduction under conditions designed to remove these three sources of optimism. \rev{Each held-out participant is excluded from z-scoring, model fitting and channel ranking.} Channel rankings are computed by cross-fitting: TreeSHAP attributions are always evaluated on participants the attributing model did not see. \rev{The primary decision rules were fixed in a statistical analysis plan before any sparse-channel result existed: the non-inferiority margin ($\delta = 0.02$ Macro-F1), the participant as the statistical unit, the bootstrap, Holm correction within each confirmatory family, and the definition of the smallest acceptable channel count $k^*$. The detailed test families for the secondary hypotheses were specified afterwards, and we report where the analysis departed from the plan (Section~\ref{sec:deviations}).} We also test a stability-regularised variant of the ranking (SR-CFA), which penalises channels whose attribution varies across disjoint groups of training participants.

The study is primarily a test, not a method proposal. We asked four prespecified questions on the 43-participant GRABMyo dataset~\cite{pradhan2022multiday}:
(H1) whether some $k < 16$ is non-inferior to all 16 channels;
(H2) whether attribution-ranked subsets beat random subsets of the same size;
(H3) whether the stability penalty improves on attribution magnitude alone;
(H4) whether the attribution is faithful, i.e.\ whether removing highly ranked channels hurts more than removing random or low-ranked ones.
We report every answer, including the negative ones. Exploratory analyses covered cross-session transfer, a bipolar re-derivation of the montage, and \rev{methodological replications on NinaPro DB5~\cite{pizzolato2017comparison} and on putEMG~\cite{kaczmarek2019putemg}, a 44-participant, 24-electrode dataset added after all other results existed}. These are reported as exploratory.
